\documentclass[11pt,a4paper]{article}

\usepackage[margin=1in]{geometry}
\usepackage{fontspec}
\usepackage{xeCJK}
\usepackage{amsmath,amssymb,amsfonts}
\usepackage{booktabs}
\usepackage{array}
\usepackage{xcolor}
\usepackage{hyperref}
\usepackage{enumitem}
\usepackage{titlesec}
\usepackage{setspace}
\usepackage{caption}
\usepackage{microtype}

\setmainfont{Times New Roman}
\setsansfont{Arial}
\setmonofont{Menlo}
\setCJKmainfont{Apple SD Gothic Neo}
\hypersetup{
  colorlinks=true,
  linkcolor=black,
  urlcolor=blue,
  citecolor=black
}

\definecolor{raiNavy}{HTML}{111827}
\definecolor{raiBlue}{HTML}{3767D8}
\definecolor{raiGreen}{HTML}{10906F}
\definecolor{raiRed}{HTML}{C0392B}
\definecolor{raiOrange}{HTML}{C45A14}
\definecolor{raiGray}{HTML}{667085}

\titleformat{\section}{\large\bfseries\color{raiNavy}}{\thesection}{0.65em}{}
\titleformat{\subsection}{\normalsize\bfseries\color{raiNavy}}{\thesubsection}{0.65em}{}
\captionsetup{font=small,labelfont=bf}
\setstretch{1.12}
\setlist[itemize]{leftmargin=1.25em,itemsep=0.25em,topsep=0.35em}

\title{\textbf{Mathematical Methodology for the RAI Risk Taxonomy Space}\\
\large Semantic State Space View of L4 Risk Cards}
\author{RAI Risk Taxonomy 2.0 Working Note}
\date{July 23, 2026}

\begin{document}
\maketitle

\begin{abstract}
This note defines the mathematical basis of the RAI Risk Taxonomy Space. The space is an exploratory semantic state space for active L4 risk cards, not a risk ranking system. Each L4 card is represented as a normalized semantic vector, projected into two dimensions for visual inspection, and interpreted relative to L3 family centroids. In a complex systems reading, L4 cards are local risk states, L3 families are semantic attractors, and HOLD cards indicate boundary regions where taxonomy assignment requires additional human review.
\end{abstract}

\section{Scope and Data Objects}

The RAI Risk Taxonomy Space visualizes the active L4 cards in release v2.17.2. The registry preserves all identifiers, while the semantic space displays only active cards.

\begin{table}[h]
\centering
\caption{Release scope used by the space view}
\begin{tabular}{llr}
\toprule
Object & Definition & Count \\
\midrule
Registered L4 IDs & Complete L4 identifier registry & 1,711 \\
Active L4 cards & Cards rendered in the semantic space & 1,660 \\
Merged records & Retired records preserved for traceability & 51 \\
HOLD cards & Active cards marked for taxonomy review & 765 \\
\bottomrule
\end{tabular}
\end{table}

Let the set of active L4 risk cards be
\begin{equation}
\mathcal{C}=\{c_i\}_{i=1}^{N}, \qquad N=1660.
\end{equation}
Each card $c_i$ contains an identifier, English and Korean labels where available, a risk definition, references, scalar risk attributes, and a released taxonomy path
\begin{equation}
\pi_i=(\ell^{(1)}_i,\ell^{(2)}_i,\ell^{(3)}_i),
\end{equation}
where $\ell^{(1)}_i$ is the L1 domain, $\ell^{(2)}_i$ is the L2 category, and $\ell^{(3)}_i$ is the L3 family.

\section{Semantic Representation}

\subsection{Card Text Construction}

For each L4 card, a semantic input string is formed from the fields that define the risk mechanism:
\begin{equation}
s_i =
\operatorname{concat}
\left(
\operatorname{label}_{i}^{en},
\operatorname{definition}_{i}^{en},
\operatorname{label}_{i}^{ko},
\operatorname{definition}_{i}^{ko}
\right).
\end{equation}
The English definition is treated as the primary semantic source. Korean text is included only as supplementary information when available. This follows the English-first display principle used in the Space interface.

\subsection{Embedding Function}

The embedding model maps each card text into a high-dimensional vector:
\begin{equation}
\mathbf{z}_i = f_{\theta}(s_i) \in \mathbb{R}^{d}.
\end{equation}
The current release uses BGE-M3 embeddings. Vectors are treated as semantic representations of failure mechanisms and evidence patterns, not as direct measurements of severity or probability.

For cosine-based comparison, vectors are normalized:
\begin{equation}
\mathbf{x}_i = \frac{\mathbf{z}_i}{\|\mathbf{z}_i\|_2}.
\end{equation}
The matrix of normalized card embeddings is
\begin{equation}
X = [\mathbf{x}_1^\top,\ldots,\mathbf{x}_N^\top]^\top \in \mathbb{R}^{N \times d}.
\end{equation}

\section{Two-Dimensional Projection}

\subsection{Deterministic PCA}

The web view uses deterministic principal component analysis (PCA) to project the active L4 embedding matrix into two dimensions. First, embeddings are centered:
\begin{equation}
\bar{\mathbf{x}} = \frac{1}{N}\sum_{i=1}^{N}\mathbf{x}_i, \qquad
\tilde{\mathbf{x}}_i=\mathbf{x}_i-\bar{\mathbf{x}}.
\end{equation}
The first two principal axes are obtained from the singular value decomposition:
\begin{equation}
\tilde{X}=U\Sigma V^\top.
\end{equation}
Let $V_2=[\mathbf{v}_1,\mathbf{v}_2]$ be the first two right singular vectors. The raw two-dimensional coordinate is
\begin{equation}
\mathbf{y}_i = \tilde{\mathbf{x}}_i V_2 \in \mathbb{R}^{2}.
\end{equation}

\subsection{Display Normalization}

To reduce the influence of extreme points in the browser view, each axis is scaled by its 99th percentile absolute coordinate:
\begin{equation}
\hat{y}_{ij} =
\operatorname{clip}\left(
\frac{y_{ij} - \bar{y}_{j}}{q_{0.99}(|y_{\cdot j}|)},
-1.2,
1.2
\right),
\qquad j \in \{1,2\}.
\end{equation}
The displayed coordinate of card $c_i$ is therefore
\begin{equation}
\hat{\mathbf{y}}_i=(\hat{y}_{i1},\hat{y}_{i2}).
\end{equation}
This scaling is purely visual. It does not change released taxonomy assignments and should not be interpreted as a calibrated risk metric.

\section{L3 Families as Semantic Attractors}

\subsection{Centroid Definition}

For each released L3 family $g \in \mathcal{G}$, define the set of active cards assigned to that family:
\begin{equation}
\mathcal{C}_g=\{c_i \in \mathcal{C}:\ell^{(3)}_i=g\}.
\end{equation}
The high-dimensional semantic centroid is
\begin{equation}
\boldsymbol{\mu}_g =
\frac{\sum_{c_i \in \mathcal{C}_g}\mathbf{x}_i}
{\left\|\sum_{c_i \in \mathcal{C}_g}\mathbf{x}_i\right\|_2}.
\end{equation}
In the browser, the visual centroid is calculated in the projected plane:
\begin{equation}
\hat{\boldsymbol{\nu}}_g =
\frac{1}{|\mathcal{C}_g|}
\sum_{c_i \in \mathcal{C}_g}\hat{\mathbf{y}}_i.
\end{equation}

\subsection{Attractor Distance}

For a selected card, the displayed distance to its released L3 family centroid is
\begin{equation}
d^{2D}_{i,g}=\left\|\hat{\mathbf{y}}_i-\hat{\boldsymbol{\nu}}_g\right\|_2.
\end{equation}
This distance is descriptive. It indicates whether the selected card is visually close to the center of its released L3 family in the two-dimensional projection. It is not a formal reassignment criterion by itself.

The high-dimensional semantic affinity used in validation is cosine similarity:
\begin{equation}
a_{i,g} = \mathbf{x}_i^\top \boldsymbol{\mu}_g.
\end{equation}
Cards with high $a_{i,g}$ are semantically close to L3 family $g$. Cards with similar affinities to multiple families are potential boundary cases.

\section{HOLD as Boundary Signal}

\subsection{Interpretation}

A HOLD marker is treated as a boundary signal, not as a separate risk type. A card can receive HOLD when the released assignment is operationally usable but still requires taxonomy review. The main sources of HOLD include:
\begin{itemize}
  \item weak separation between competing L3 families,
  \item L4 labels or definitions that are noisy or insufficiently specific,
  \item evidence sources that support a broad topic but not the specific L4 failure mechanism,
  \item missing or under-defined L3 families in the current taxonomy.
\end{itemize}

\subsection{Quantitative Boundary Form}

Let $g_1(i)$ and $g_2(i)$ denote the two L3 families with the highest semantic affinities for card $c_i$:
\begin{equation}
g_1(i)=\arg\max_{g\in\mathcal{G}} a_{i,g},
\end{equation}
\begin{equation}
g_2(i)=\arg\max_{g\in\mathcal{G}\setminus\{g_1(i)\}} a_{i,g}.
\end{equation}
The semantic margin is
\begin{equation}
m_i=a_{i,g_1(i)}-a_{i,g_2(i)}.
\end{equation}
A small $m_i$ indicates that the card is close to a boundary between L3 attractors. A review flag can therefore be expressed as
\begin{equation}
H_i =
\mathbb{I}\left[
m_i < \tau_m
\right]
\lor
\mathbb{I}\left[
a_{i,\ell^{(3)}_i}<\tau_a
\right]
\lor
Q_i,
\end{equation}
where $\tau_m$ is a margin threshold, $\tau_a$ is an assignment-affinity threshold, and $Q_i$ is a qualitative review indicator for label-definition mismatch, weak evidence, or unresolved taxonomy coverage. The public Space displays $H_i$ as a compact HOLD marker.

\section{Complex Systems Reading}

\subsection{State Space}

The Space can be read as a semantic state space:
\begin{equation}
\mathcal{S} = \left(\mathcal{C}, X, \mathcal{G}, \Pi, H\right),
\end{equation}
where $\mathcal{C}$ is the set of L4 states, $X$ is the embedding matrix, $\mathcal{G}$ is the set of L3 families, $\Pi$ is the released assignment function, and $H$ is the HOLD review indicator.

In this interpretation:
\begin{itemize}
  \item L4 cards are local risk states.
  \item L3 families are semantic attractors that gather similar failure mechanisms.
  \item HOLD cards are boundary states where assignment is less stable.
  \item Physical AI forms a comparatively stable subsystem because many of its risk mechanisms involve concrete sensing, embodiment, and physical intervention.
\end{itemize}

\subsection{Interface Modes}

The Space provides three visual interpretation modes.

\begin{table}[h]
\centering
\caption{Interpretation modes in the Space interface}
\begin{tabular}{p{0.22\linewidth}p{0.68\linewidth}}
\toprule
Mode & Interpretation \\
\midrule
Domain fields & Colors encode L1 domains: General-purpose AI, Agentic AI, and Physical AI. \\
Boundary regions & HOLD cards are emphasized to reveal regions of taxonomy uncertainty or category overlap. \\
L3 attractors & L3 centroids are shown to inspect whether selected cards are central or peripheral within their released L3 family. \\
\bottomrule
\end{tabular}
\end{table}

\section{Reliability and Limits}

The Space supports taxonomy diagnosis, not automatic ground-truth revision. The following limits apply:
\begin{itemize}
  \item PCA preserves major variance directions but does not preserve all high-dimensional neighborhoods.
  \item A short 2D distance does not prove that two L4 cards are duplicates.
  \item A long 2D distance does not prove that a released L3 assignment is wrong.
  \item HOLD should be interpreted jointly with semantic margins, card definitions, and evidence quality.
  \item Human review remains necessary for taxonomy changes, especially when a card spans multiple L3 families.
\end{itemize}

The practical role of the Space is therefore diagnostic. It helps identify crowded attractors, unstable boundaries, under-specified L3 definitions, and candidate regions where future taxonomy refinement may improve reliability.

\section{Summary}

The mathematical pipeline is:
\begin{equation}
\text{L4 card text}
\rightarrow
\text{BGE-M3 embedding}
\rightarrow
\text{PCA projection}
\rightarrow
\text{L3 centroid comparison}
\rightarrow
\text{complex systems interpretation}.
\end{equation}
Under this view, the RAI Risk Taxonomy Space is a semantic state-space instrument for inspecting how active L4 risk cards cluster around L3 attractors and where taxonomy boundary conditions require review.

\end{document}
